\section{The sum connection on both half-line branches and its exact domain}
\label{sec:sum-connection-stieltjes}

This section integrates the complete sum-connection delivery, whose
source note is retained in the appendices at its declared PR~20 revision
\nolinkurl{2144c358c2b41aa235b15cf0fa472289aec8673f}. We prove its
mass-dependent derivative identity in the original source and then carry
the full relative connection into the quadratic coordinate. Both branches,
all relative coefficients, and the condition where the two branches meet
are retained. Only the stated dagger stability is needed until evenness
is explicitly imposed.

\subsection{The actual amplitude, its mass, and the original relative fibre}

Keep \(g=2\xi\), \(v_h=g/h\), the entire quotient at the full selected
zero orders, and the original \(F_h\in\mathscr B\) with Mellin transform
\(v_h\). Let the packet be stable under \(\rho\mapsto1-\overline\rho\).
The empty polynomial \(h=1\) is also allowed for this analytic calculation;
no nonzero finite jet quotient is inferred in that case. Set
\[
 a_h^{\rm SC}(t)=\frac{v_h(1/2+it)}{\sqrt{2\pi}},\quad
 w_h(t)=|a_h^{\rm SC}(t)|^2,\quad
 \mu_h=\int_{\mathbb R}w_h(t)dt=\|F_h\|_{L^2(dx)}^2.
 \tag{FC.1}
\]
The superscript distinguishes the delivered spectral amplitude from
the gamma-reference multiplier in (TG.4). Their exact relation is
\(a_h^{\rm SC}(t)=\Gamma(1/4+it/2)a_h^{\rm TG}(t)/\sqrt{2\pi}\).
Taking its squared modulus gives precisely the two original measure
formulas; neither mass nor phase is replaced.

For \(k\ge2\) use
\[
 u=\sum_i t_i,\quad y_i=t_i-u/k\ (i<k),\quad
 y_k=-\sum_{i<k}y_i,\quad t_i=u/k+y_i,
 \qquad \Psi(u,\mathbf y)=\prod_i a_h^{\rm SC}(t_i).
 \tag{FC.2}
\]
The Jacobian has absolute value one. One way to check its constant
is to use coordinates \((u,t_1,\ldots,t_{k-1})\), whose inverse
sets \(t_k=u-\sum_{i<k}t_i\) and has determinant of absolute
value one. The replacement \(t_i=y_i+u/k\), \(i<k\), is triangular
with diagonal one. For \(k=2\), \(v=t_1-t_2=2y_1\), so the
earlier measure \(du\,dv/2\) is unchanged.

The test-space estimates imply that \(e^{r/2}F_h(e^r)\) and all
its derivatives and polynomial multiples are Schwartz on \(\mathbb R\).
Indeed arbitrary powers of \(e^r\) bound both tails, and the
logarithmic derivatives are the stated \(D\)-derivatives of \(F_h\).
Fourier differentiation and integration by parts therefore show that
\(a_h^{\rm SC}\) and all its derivatives are Schwartz. These facts
justify the product integrals and differentiated integrals below.

Dagger stability gives \(\overline{a_h^{\rm SC}}=(-1)^d a_h^{\rm SC}\).
Thus this amplitude and its derivative have the same constant real or
imaginary phase. Off the discrete zero set,
\(w_h'=2\overline{a_h^{\rm SC}}(a_h^{\rm SC})'\), with the
product real, and consequently
\[
 I_h:=\int\frac{|w_h'|^2}{w_h}
   =4\int|(a_h^{\rm SC})'|^2
   =4\int_0^\infty(\log x)^2|F_h(x)|^2dx,
 \qquad \int\overline{a_h^{\rm SC}}(a_h^{\rm SC})'=0.
 \tag{FC.3}
\]
Values assigned to the quotient on its discrete zero set do not
change these integrals. The second equality is the original Mellin
Plancherel identity; the last integral is \(\frac12\int w_h'=0\)
by Schwartz decay. In particular every number in (FC.3) is finite.

Put \(m(u)=\|\Psi(u,\cdot)\|^2=w_h^{*k}(u)>0\). Positivity holds
at each \(u\): the factors have only discrete zeros, whose inverse
images in a relative fibre are a countable union of null hyperplanes.
The integral of \(m\) is exactly \(\mu_h^k\). Let \(\Pi_{0,u}\)
be orthogonal projection onto the line spanned by \(\Psi(u,\cdot)\),
and let \(n(u)=(1-\Pi_{0,u})\partial_u\Psi\). Constant phase gives
\(\langle\Psi,\partial_u\Psi\rangle=m'/2\), so
\[
 n=\partial_u\Psi-\frac{m'}{2m}\Psi,\qquad
 \boxed{\int\frac{|m'|^2}{m}+4\int\|n\|^2
                   =\frac{\mu_h^{k-1}}k I_h.}
 \tag{FC.4}
\]
To prove the equality, differentiate the product in (FC.2).
Its derivative is
\[
 k^{-1}\sum_i(a_h^{\rm SC})'(t_i)
                \prod_{j\ne i}a_h^{\rm SC}(t_j).
\]
Before the common factor \(k^{-2}\), each of the \(k\) diagonal
terms in its squared norm has integral
\(\mu_h^{k-1}\|(a_h^{\rm SC})'\|^2\), and every distinct-index
term contains \(|\int\overline{a_h^{\rm SC}}(a_h^{\rm SC})'|^2=0\).
Thus its integral is \(\mu_h^{k-1}I_h/(4k)\). The orthogonal
line and complementary components have squared norms
\(|m'|^2/(4m)\) and \(\|n\|^2\); integration proves (FC.4).

\subsection{The complete matrix connection and its actual graph}

Choose a fixed finite independent homogeneous relative polynomial
family \(\theta_\alpha\), of degrees \(\ell_\alpha\), in the
\(k-1\) independent variables \(z_i=s_i-S/k\). At source degree
\(M\), retain exactly
\[
 P(\mathbf s)=\sum_\alpha\theta_\alpha(\mathbf z)f_\alpha(S),
 \quad \deg f_\alpha\le M-\ell_\alpha,
 \quad S=k/2+iu,
 \quad c_P(u)=(f_\alpha(k/2+iu))_\alpha.
 \tag{FC.5}
\]
The family can be the whole admitted relative monomial family or the
specified invariant family. Its cardinality \(r\) and all degree
bounds remain part of the source; it is not held fixed while \(M\)
changes. In \(\mathscr H=L^2(\mathbb R^{k-1},d\mathbf y)\), define
\[
 \begin{gathered}
 j_uc=\Psi(u,\mathbf y)\sum_\alpha\theta_\alpha(i\mathbf y)c_\alpha,
 \quad W=j_u^*j_u>0,\quad B=j_u^*j_u',\quad T=(j_u')^*j_u',\\
 \Pi=j_uW^{-1}j_u^*,\quad \Gamma=W^{-1}B,\quad
 N=(1-\Pi)j_u',\quad \mathcal N=N^*N=T-B^*W^{-1}B.
 \end{gathered}
 \tag{FC.6}
\]
Positivity follows from independence of the relative polynomials and
the almost-everywhere nonvanishing just used. Schwartz estimates on
compact \(u\)-intervals prove smoothness of these finite matrices
and of \(j_u\) in operator norm. Direct differentiation gives
\[
 \begin{gathered}
 \partial_u(j_uc)=j_u(c'+\Gamma c)+Nc,\qquad j_u^*N=0,\\
 \Gamma^*W+W\Gamma=W',\qquad \mathcal N\succeq0,\\
 \|\partial_u(j_uc)\|^2
       =(c'+\Gamma c)^*W(c'+\Gamma c)+c^*\mathcal Nc.
 \end{gathered}
 \tag{FC.7}
\]
For example \(j_u\Gamma=\Pi j_u'\); adding the complementary
term proves the first identity. Differentiating \(W=j_u^*j_u\)
proves metric compatibility, and the projection identity proves the
normal Gram and positivity. In this original fixed frame the common
phase makes \(B=B^*\), hence \(B=W'/2\). After a variable change
of frame \(j^C=jC(u)\), one instead keeps the full formulas
\[
 \begin{gathered}
 W^C=C^*WC,\quad B^C=C^*BC+C^*WC',\\
 \Gamma^C=C^{-1}\Gamma C+C^{-1}C',\quad N^C=NC,
 \quad \mathcal N^C=C^*\mathcal NC.
 \end{gathered}
 \tag{FC.8}
\]
They follow by the product rule and the unchanged projection image.
For instance \(C=e^{iu}I\) gives \(B^C=B+iW\), so the fixed-frame
identity \(B=W'/2\) is not silently imposed on a changed frame.

The actual derivative graph is the image of
\(P\mapsto(j_u(c_P'+\Gamma c_P),Nc_P)\). Addition is an
isometric injection on that graph by (FC.7); its image is precisely
\(\partial_u\mathscr U_k\mathcal T_h^{(k)}P\). Here \(\mathscr U_k\)
is the original Mellin--Fourier isometry followed by (FC.2). On the
test space it obeys
\[
 \mathscr U_kD^{(k)}=(k/2+iu)\mathscr U_k,
 \quad \partial_u\mathscr U_k=i\mathscr U_k\mathscr L_k,
 \quad \mathscr L_k=\frac1k\sum_i\log x_i,
 \quad [D^{(k)},\mathscr L_k]=-I.
 \tag{FC.9}
\]
The first two formulas follow by integration by parts and
differentiation of the Mellin integral. Applying
\(-x_i\partial_{x_i}\) to \(\log x_i\) gives the last formula,
including \(1/k\). Therefore addition of both graph components,
followed by \(-i\mathscr U_k^{-1}\) and the original full jet,
recovers \(J^{(k)}\mathscr L_k\mathcal T_h^{(k)}P\). The bare
derivative lands in the image of graph addition; to land in the
pair graph itself one uses the actual map
\(v\mapsto(\Pi_u\partial_uv,(1-\Pi_u)\partial_uv)\).

\subsection{A branch-preserving quadratic map for the entire matrix weight}

Put \(x=u^2>0\), \(v=\sqrt x\), and retain the ordered branches
\(u=v\) and \(u=-v\). These letters denote the spectral sum
coordinate, not the original positive Mellin variables in (FC.9).
For any coefficient function define
\[
 c(v)=a(x)+ivb(x),\quad c(-v)=a(x)-ivb(x),\quad
 \mathsf M(x)=\begin{pmatrix}I&ivI\\I&-ivI\end{pmatrix}.
 \tag{FC.10}
\]
Its inverse is \(a=(c(v)+c(-v))/2\),
\(b=(c(v)-c(-v))/(2iv)\). No relation between \(W(v)\) and
\(W(-v)\) is assumed. Write \(W_\pm=W(\pm v)\). The full
positive paired density is
\[
 \begin{aligned}
 \mathsf H(x)&=\mathsf M^*
          \frac{\operatorname{diag}(W_+,W_-)}{2v}\mathsf M\\
 &=\frac1{2\sqrt x}
 \begin{pmatrix}
 W_++W_-&i\sqrt x(W_+-W_-)\\
 -i\sqrt x(W_+-W_-)&x(W_++W_-)
 \end{pmatrix}>0.
 \end{aligned}
 \tag{FC.11}
\]
The inverse and the change of variable on each half-line prove the
unitary equality
\[
 \int_{\mathbb R}c(u)^*W(u)c(u)du
   =\int_0^\infty \binom ab^*\mathsf H(x)\binom ab\,dx.
 \tag{FC.12}
\]
Here \(\binom ab^*\) means the adjoint of the two-block column.
Positivity follows from the invertibility of \(\mathsf M\) at
every \(x>0\); surjectivity on these \(L^2\) spaces follows
from its pointwise inverse and the same norm identity. Both branch
weights and their off-diagonal difference remain in (FC.11).

For the precise polynomial image in (FC.5), write
\(f_\alpha(S)=\sum_n f_{\alpha,n}S^n\). Then
\[
 \begin{aligned}
 a_\alpha(x)&=\sum_n f_{\alpha,n}\sum_{2j\le n}
           \binom n{2j}(k/2)^{n-2j}(-x)^j,\\
 b_\alpha(x)&=\sum_n f_{\alpha,n}\sum_{2j+1\le n}
           \binom n{2j+1}(k/2)^{n-2j-1}(-x)^j.
 \end{aligned}
 \tag{FC.13}
\]
Their degree bounds are respectively
\(\lfloor(M-\ell_\alpha)/2\rfloor\) and
\(\lfloor(M-\ell_\alpha-1)/2\rfloor\), with a negative bound
meaning a zero component. The inverse is exactly
\(f_\alpha(S)=a_\alpha(-(S-k/2)^2)
 +(S-k/2)b_\alpha(-(S-k/2)^2)\). Expanding the original powers
proves both formulas and their inverse, including every \(k/2\)
coefficient. Thus (FC.12) retains the full filtered source.

\begin{theorem}[The complete paired connection and normal energy]
Define
\[
 \begin{gathered}
 \Gamma_b(x)=\operatorname{diag}
       \left(\frac{\Gamma(v)}{2v},-\frac{\Gamma(-v)}{2v}\right),\\
 \mathsf A(x)=\mathsf M^{-1}\Gamma_b\mathsf M
                 +\mathsf M^{-1}\mathsf M'
       =\mathsf M^{-1}\Gamma_b\mathsf M
                 +\operatorname{diag}(0,(2x)^{-1}I),\\
 \mathsf Q(x)=\mathsf M^*
       \frac{\operatorname{diag}(\mathcal N(v),\mathcal N(-v))}{2v}
                                    \mathsf M\succeq0,
 \qquad \nabla_x=\partial_x+\mathsf A.
 \end{gathered}
 \tag{FC.14}
\]
For \(q=(a,b)^T\), the exact differentiated norm is
\[
 \boxed{\int_{\mathbb R}\|\partial_u(j_uc(u))\|^2du
       =4\int_0^\infty x(\nabla_xq)^*\mathsf H\nabla_xq\,dx
                 +\int_0^\infty q^*\mathsf Qq\,dx.}
 \tag{FC.15}
\]
The transformed coefficient derivative and its metric identities are
\[
 \begin{gathered}
 \mathsf J(x)\nabla_x,
 \quad \mathsf J=\begin{pmatrix}0&2ixI\\-2iI&0\end{pmatrix},
 \quad \mathsf J^2=4xI,
 \quad \mathsf J^*\mathsf H\mathsf J=4x\mathsf H,\\
 \mathsf A^*\mathsf H+\mathsf H\mathsf A
          =\mathsf H'+\frac1{2x}\mathsf H.
 \end{gathered}
 \tag{FC.16}
\]
Consequently \(\mathsf A-(4x)^{-1}I\) is the connection matrix
compatible with the full density \(\mathsf H\). The explicit
Jacobian term in (FC.16) remains in the transported original derivative.
\end{theorem}

\begin{proof}
Let \(F=(c(v),c(-v))^T=\mathsf Mq\) and
\(\mathsf K=\operatorname{diag}(2vI,-2vI)\). The two original
covariant derivative values are exactly
\(\mathsf K(F'+\Gamma_bF)=\mathsf K\mathsf M\nabla_xq\).
The minus sign on the second branch is from \(d(-v)/dx=-1/(2v)\).
Multiplying the explicit inverse of \(\mathsf M\) gives
\(\mathsf M^{-1}\mathsf K\mathsf M=\mathsf J\) and
\(\mathsf M^{-1}\mathsf M'=\operatorname{diag}(0,1/(2x)I)\).
The block-scalar matrix \(\mathsf K\) commutes with the two-block
weight and its squared modulus is \(4xI\). This proves the first
energy term and the first line of (FC.16). Applying (FC.12) to
the original normal Gram proves the second energy term; (FC.7)
then proves (FC.15).

For metric compatibility set
\(\mathsf D=\operatorname{diag}(W_+,W_-)/(2v)\).
The original equality \(\Gamma^*W+W\Gamma=W'\), with the
two opposite derivatives of \(v\), gives
\(\Gamma_b^*\mathsf D+\mathsf D\Gamma_b
 =\mathsf D'+\mathsf D/(2x)\). The additional term is the
derivative of the factor \(1/(2v)\). Substitution of
\(\mathsf H=\mathsf M^*\mathsf D\mathsf M\) and the two
terms defining \(\mathsf A\) proves the last line of (FC.16).
Subtracting \(1/(4x)I\) from the matrix and its adjoint removes
exactly that extra term, proving the last assertion.
\end{proof}

Multiplication by the original spectral sum has the exact paired matrix
\[
 \mathsf S=k/2\,I+
                \begin{pmatrix}0&-xI\\I&0\end{pmatrix}.
 \tag{FC.17}
\]
This follows by multiplying \(a+iu b\) by \(k/2+iu\).
On smooth compactly supported sections in \(x>0\), the commutator
is \([\mathsf J\nabla_x,\mathsf S]=iI\). To prove it with
all variable-frame terms present, conjugate to the two branch
coordinates used in the proof: multiplication is
\(\operatorname{diag}(k/2+iv,k/2-iv)\), the connection part
commutes with this block-scalar matrix, and multiplication of its
derivative by \(\mathsf K\) gives \(iI\) on both blocks.
Conjugating back proves the claim. Through (FC.9) this is the
same original pair \(D^{(k)},\mathscr L_k\), with its original sign.

\subsection{The meeting point of the two branches is part of the domain}

For the \(L^2\) identity (FC.12) no value at zero is required.
For the derivative domain it is essential. Define the original domain
to consist of coefficient functions \(c\) that are locally absolutely
continuous on \(\mathbb R\), have finite norm in (FC.12), and have
finite derivative energy in (FC.7). Equivalently these are the
coefficient functions for which \(j_uc(u)\) has an \(L^2\) weak
derivative: on each compact interval the smooth left inverse
\(W^{-1}j_u^*\) and strict positivity of \(W\) prove both directions
by the product rule. Its exact image under (FC.10) consists of
locally absolutely continuous \(q\) on \((0,\infty)\) with the two
finite integrals in (FC.12), (FC.15), and the joining condition
\[
 \lim_{x\downarrow0}(a(x)+i\sqrt x\,b(x))
     =\lim_{x\downarrow0}(a(x)-i\sqrt x\,b(x)).
 \tag{FC.18}
\]
These finite one-sided limits exist under the stated integrability
conditions. Indeed \(W\) is bounded above and below by positive
constants near zero and \(\Gamma\) is bounded there. The two
finite integrals imply that each branch coefficient and its ordinary
\(u\)-derivative are square-integrable near zero. The fundamental
theorem for locally absolutely continuous functions and
Cauchy--Schwarz then give their finite endpoint limits. Their union
has a weak derivative in \(L^2\) precisely when those limits agree;
integration by parts against a compactly supported test function
shows that a mismatch contributes their nonzero jump times a delta
distribution. This proves (FC.18) and both directions of the
domain correspondence. For the original polynomial image (FC.13),
the condition holds automatically. In general, finite energy on the
two open branches alone does not impose it: a function constant with
different values on the two sides near zero has zero ordinary
derivative there, while its joined distributional derivative has
that nonzero jump. The matching condition preserves the original
domain without discarding either branch.

\subsection{The even specialization and the unchanged scalar contraction}

If the original packet also has conjugation closure, \(w_h\) is even.
For the symmetric square with relative frame \(1,\Delta,\Delta^2,\ldots\),
the actual \(j_u\) is even in \(u\), since its two amplitude
factors are interchanged by \(u\mapsto-u\) and amplitude evenness.
Thus \(W\) and \(\mathcal N\) are even and \(\Gamma\) is odd.
Set \(L(x)=W(\sqrt x)/\sqrt x\),
\(L_N(x)=\mathcal N(\sqrt x)/\sqrt x\), and
\(\Gamma_x=\Gamma(\sqrt x)/(2\sqrt x)\). Equations (FC.14)--(FC.16)
reduce, by direct block substitution, to
\[
 \begin{gathered}
 \mathsf H=\operatorname{diag}(L,xL),\qquad
 \mathsf Q=\operatorname{diag}(L_N,xL_N),\qquad
 \mathsf A=\operatorname{diag}(\Gamma_x,\Gamma_x+(2x)^{-1}I),\\
 \int\|\partial_u(j_uc)\|^2du
 =\int_0^\infty\bigl[
 4x(a'+\Gamma_xa)^*L(a'+\Gamma_xa)
 +(b+2x(b'+\Gamma_xb))^*L(b+2x(b'+\Gamma_xb))\\
 \hspace{35mm}{}+a^*L_Na+x b^*L_Nb\bigr]dx.
 \end{gathered}
 \tag{FC.19}
\]
These are exactly the two measures and the complete relative family
of the preceding symmetric-square section. Its unequal finite
quotient components are obtained from this same filtered source,
using its already proved full jet and actual minimum.

For any \(k\ge2\) under the same quartet assumption, the scalar sum
density \(m\) is even. Define the actual half-line density
\(\rho_k(x)=m(\sqrt x)/\sqrt x\). It retains total mass
\(\int_0^\infty\rho_k=\mu_h^k\). The source contraction (FC.4)
has the precise radial form
\[
 \boxed{
 4\int_0^\infty x\,
       \frac{|\rho_k'(x)+\rho_k(x)/(2x)|^2}{\rho_k(x)}dx
 +4\int_0^\infty\frac{\|n(\sqrt x)\|^2}{\sqrt x}dx
                =\frac{\mu_h^{k-1}}k I_h.}
 \tag{FC.20}
\]
To verify it, \(m(u)=u\rho_k(u^2)\) for \(u>0\), so
\(m'(u)=\rho_k(x)+2x\rho_k'(x)\). Changing variables in
\(2\int_0^\infty m'^2/m\,du\) gives its first term.
Reflection of all relative variables is a unitary map of each
fibre and carries \(n(u)\) to \(-n(-u)\), so its norm is even;
changing variables gives the second term. This proves (FC.20)
with the stated drift and the entire relative contribution.
It is the transported source estimate, with its mass and domain.
The complete matrix energy remains (FC.15), on the actual graph
and with its normal Gram, rather than being assigned the scalar
constant-column bound.

All linear maps on the displayed filtered sources and graph images
have the original fixed-support lift: they preserve the support
label, send a supported zero to the corresponding supported zero,
and send external absence to external absence. Their additive and
scalar identities follow from the proved linear maps. In particular
the arithmetic observation is taken after graph addition and the
specified inverse Mellin isometry; no arithmetic jet is assigned
to an arbitrary independent measurable normal section.

\subsection{The actual arithmetic normal contribution is strictly positive}

We now use the actual zero set of the amplitude in (FC.1).
Hardy's theorem that infinitely many zeros of the Riemann zeta
function lie on the critical line, with the original source and
reading record retained in Section~\ref{sec:theta-gamma-reference},
implies that
\[
 \mathcal Z_h=\{T\in\mathbb R:a_h^{\rm SC}(T)=0\}
 \quad\hbox{is infinite and discrete.}
 \tag{FC.21}
\]
Indeed a finite polynomial \(h\) removes only its selected finite
set of complete zero orders. Every remaining critical-line zero is
still a zero of \(v_h\), and the nonzero entire function \(v_h\)
has a discrete zero set. This also applies to \(h=1\).
At any \(T\in\mathcal Z_h\), keep its entire order
\(\ell_T\ge1\) and its original leading coefficient
\[
 a_h^{\rm SC}(T+\delta)
       =c_T\delta^{\ell_T}+O(\delta^{\ell_T+1}),\qquad
 (a_h^{\rm SC})'(T+\delta)
       =\ell_Tc_T\delta^{\ell_T-1}+O(\delta^{\ell_T}),
 \quad c_T=\frac{(a_h^{\rm SC})^{(\ell_T)}(T)}{\ell_T!}\ne0.
 \tag{FC.22}
\]
These are expansions of the original amplitude, including its fixed
phase and its factor \(1/\sqrt{2\pi}\).

\begin{theorem}[Strict scalar residual and every finite normal direction]
For the actual amplitude and every integer \(k\ge2\),
\[
 \int_{\mathbb R}\|n(u)\|^2du>0,\qquad
 \boxed{\int_{\mathbb R}\frac{|m'(u)|^2}{m(u)}du
                <\frac{\mu_h^{k-1}}k I_h.}
 \tag{FC.23}
\]
For \(k\ge3\), \(\|n(u)\|>0\) for every real \(u\).
For any nonempty finite independent relative polynomial frame of
(FC.5), its full normal Gram satisfies
\[
 \boxed{\begin{aligned}
  \mathcal N(u)&>0&&\text{for every }u\in\mathbb R,
                                      &&k\ge3,\\
  \mathcal N(u)&>0&&\text{for every }u\notin
                  \mathcal Z_h+\mathcal Z_h,
                                      &&k=2.
 \end{aligned}}
 \tag{FC.24}
\]
In particular \(\mathcal N(u)>0\) almost everywhere for every
\(k\ge2\), since \(\mathcal Z_h+\mathcal Z_h\) is countable.
No lower bound uniform in \(u\), the relative degree, or the
packet is asserted by these strict inequalities.
\end{theorem}

\begin{proof}
First fix \(T\in\mathcal Z_h\) and choose
\(b_2,\ldots,b_k\) outside \(\mathcal Z_h\). Put
\(u_0=T+\sum_{i=2}^kb_i\) and
\(A_b=\prod_{i=2}^ka_h^{\rm SC}(b_i)\ne0\).
Within that fixed original sum fibre consider the transverse path
\[
 t_1=T+\delta,\quad t_i=b_i\ (2\le i<k),\quad
 t_k=b_k-\delta.
\]
It changes the relative variables while retaining \(u=u_0\).
The derivative \(\partial_u\) is still the original derivative
at fixed relative variables in (FC.2), and its product formula gives
\[
 \begin{aligned}
  \Psi(u_0,\mathbf y(\delta))
      &=c_TA_b\delta^{\ell_T}+O(\delta^{\ell_T+1}),\\
  \partial_u\Psi(u_0,\mathbf y(\delta))
      &=\frac{\ell_T}{k}c_TA_b\delta^{\ell_T-1}
                              +O(\delta^{\ell_T}).
 \end{aligned}
 \tag{FC.25}
\]
Only the derivative of the first amplitude has order
\(\ell_T-1\); every other product term retains the full
factor \(a_h^{\rm SC}(T+\delta)\) and has order at least
\(\ell_T\). Since \(m(u_0)>0\) and \(m\) is smooth,
subtracting \(m'(u_0)\Psi/(2m(u_0))\) does not change that
nonzero leading coefficient. Thus \(n(u_0,\cdot)\) is nonzero
at nearby points of this path and, by continuity, on a nonempty
open subset of the entire relative fibre. Its squared fibre norm
is strictly positive. Differentiated Schwartz estimates show that
\(\Psi\) and \(\partial_u\Psi\) are continuous as functions
of \(u\) with values in \(\mathscr H\). Smoothness and positivity
of \(m\) therefore make \(n\) continuous in that Hilbert norm.
Its squared norm stays positive on some interval around \(u_0\),
proving positive integrated residual. Equation (FC.4) proves (FC.23).

For \(k\ge3\), the same construction works at every prescribed
\(u_0\). Choose \(b_3,\ldots,b_{k-1}\) outside the zero set
if that list is nonempty. Choose \(b_2\) outside both
\(\mathcal Z_h\) and the translate-reflection of that set which
would make
\(b_k=u_0-T-\sum_{i=2}^{k-1}b_i\) a zero.
These two discrete sets cannot exhaust the real line. This proves
the pointwise scalar statement.

To prove the full matrix assertion, fix \(u\) and suppose
\(N_uc=0\). The exact projection formula gives
\(j_u'c=j_ud\), where \(d=\Gamma_uc\) is a finite coefficient
vector. Write \(\vartheta_c(\mathbf y)=
\sum_\alpha\theta_\alpha(i\mathbf y)c_\alpha\), and similarly
\(\vartheta_d\). Equality in \(\mathscr H\) says
\[
 (\partial_u\Psi)\vartheta_c=\Psi\vartheta_d.
\]
Both sides have smooth representatives in the relative variables,
so their almost-everywhere equality is equality everywhere: a nonzero
continuous difference at a point would remain nonzero on a set of
positive measure.

Suppose \(k\ge3\). For each \(T\in\mathcal Z_h\), restrict
to the affine hyperplane \(t_1=u/k+y_1=T\), of dimension
\(k-2\ge1\). Each of the other original coordinates is a
nonconstant affine function on this hyperplane. The points where
one of their amplitude factors vanishes are consequently a countable
union of proper affine hyperplanes of measure zero. Their complement
is dense and is open near each of its points. At any such point
\(p\), apply (FC.25). In the asserted smooth identity the left
side has leading coefficient
\(\ell_Tc_TA_b\vartheta_c(p)/k\) in degree \(\ell_T-1\),
whereas the right side has order at least \(\ell_T\).
Hence \(\vartheta_c(p)=0\).
Continuity proves its vanishing on the entire real hyperplane.
Division in \(y_1\) then shows that the polynomial
\(\vartheta_c\) has factor \(y_1-(T-u/k)\): the remainder
polynomial in the other variables vanishes on all their real values,
and therefore is identically zero.
There are infinitely many distinct such \(T\), while the degree
in \(y_1\) is finite. Thus \(\vartheta_c=0\), and independence
of the original relative polynomials gives \(c=0\).
This proves injectivity of \(N_u\) and therefore
\(\mathcal N(u)>0\) on the whole finite coefficient space.

For \(k=2\), take \(u\notin\mathcal Z_h+\mathcal Z_h\).
At every \(T\in\mathcal Z_h\), the other amplitude
\(a_h^{\rm SC}(u-T)\) is nonzero. The same expansion forces
the univariate polynomial \(\vartheta_c\) to vanish at
\(y_1=T-u/2\). These are infinitely many distinct points,
so again \(\vartheta_c=0\) and \(c=0\).
The discrete subset \(\mathcal Z_h\) of \(\mathbb R\)
is countable, hence so is its set of pairwise sums. This proves
every assertion in (FC.24).
\end{proof}

The consequences hold for the actual full derivative graph and both
quadratic branches. For every nonzero admitted source polynomial
\(P\),
\[
 \begin{gathered}
 0<\int_{\mathbb R}c_P(u)^*\mathcal N(u)c_P(u)\,du<\infty,\\
 \mathsf Q(x)>0\quad\text{for all }x>0\text{ when }k\ge3,
 \qquad
 \mathsf Q(x)>0\quad\text{for almost every }x>0\text{ when }k=2.
 \end{gathered}
 \tag{FC.26}
\]
Indeed a nonzero vector polynomial \(c_P\) is nonzero outside
a finite set on the real line, and (FC.24) supplies strict
positivity almost everywhere. The integral is finite because
\(\|N_uc_P\|\le\|j_u'c_P\|\), and the latter is a
polynomial multiple of differentiated Schwartz columns.
The assertion about \(\mathsf Q\) follows from its invertible
congruence (FC.14). For \(k=2\), the exceptional positive
\(x\) for which either \(\sqrt x\) or \(-\sqrt x\) belongs
to \(\mathcal Z_h+\mathcal Z_h\) form a countable set.
Under every declared smooth invertible frame,
\(N_u^Cc_P^C(u)=N_uc_P(u)\) and
\(\mathcal N^C=C^*\mathcal NC\).
Thus the full normal energy and strictness are preserved by the
exact coefficient isomorphism \(c_P^C=C^{-1}c_P\), including
its original filtered image.

The excluded two-factor fibre can actually vanish. For the quartet
packet, the original amplitude is even, so
\[
 \partial_u\Psi(0,y)=\tfrac12\bigl(
 (a_h^{\rm SC})'(y)a_h^{\rm SC}(-y)
 +a_h^{\rm SC}(y)(a_h^{\rm SC})'(-y)\bigr)=0.
 \tag{FC.27}
\]
Since the relative frame is independent of \(u\), this gives
\(j_0'=0\), \(N_0=0\), and \(\mathcal N(0)=0\).
Also \(m'(0)=0\), hence \(n(0)=0\) exactly. The positive
total residual and the almost-everywhere full matrix assertion
therefore retain this original zero fibre. They prove that every
nonzero actual finite polynomial direction has a positive normal
contribution while leaving its complete size in (FC.15) and (FC.26).

\subsection{Exact reconstruction from the actual finite normal image}

Fix \(k\ge2\), the same actual packet, and one of the finite
relative families in (FC.5). Retain its original degree bound \(M\)
and the literal polynomial basis
\[
 \begin{gathered}
 \mathcal I_M=\{(\alpha,n):0\le n\le M-\ell_\alpha\},
 \quad q_M=|\mathcal I_M|,\qquad
 P_{\alpha,n}(\mathbf s)=\theta_\alpha(\mathbf z)S^n,\\
 P_p=\sum_{(\alpha,n)\in\mathcal I_M}p_{\alpha,n}P_{\alpha,n},
 \quad c_{P_p}(u)=C_M(u)p,
 \quad (C_M(u))_{\beta,(\alpha,n)}
       =\delta_{\alpha\beta}(k/2+iu)^n.
 \end{gathered}
 \tag{FC.28}
\]
Thus \(C_M\) is a specified rectangular polynomial matrix, not a
change of relative frame. Each of its columns retains the original
power of \(S\). Its derivative is
\((C_M')_{\beta,(\alpha,n)}=
\delta_{\alpha\beta}in(k/2+iu)^{n-1}\) for \(n\ge1\),
and zero for \(n=0\). Independence of the relative family over
\(\mathbb C\), followed by comparison of the powers of \(S\),
proves uniqueness of these coefficients. The source is the exact
space \(\mathcal P_M^\theta=\operatorname{span}\{P_{\alpha,n}\}\);
it is the full original admitted source when that family is the
complete admitted relative basis, and the declared summand otherwise.
If \(\mathcal I_M\) is empty, all the finite source and image
spaces below are zero and their maps are the unique maps of zero
spaces. We write the nonempty case explicitly.

Let \(\mathscr K=L^2(\mathbb R_u;\mathscr H)\), with the
original measures \(du\,d\mathbf y\). Define the normal and
tangential maps by their actual amplitudes,
\[
 \begin{aligned}
 \mathbf N_M:\mathbb C^{q_M}&\longrightarrow\mathscr K,
       & (\mathbf N_Mp)(u)&=N_uC_M(u)p,\\
 \mathbf T_M:\mathbb C^{q_M}&\longrightarrow\mathscr K,
       & (\mathbf T_Mp)(u)&=j_u(C_M'(u)+\Gamma_uC_M(u))p,\\
 H_M^N&=\mathbf N_M^*\mathbf N_M
       =\int_{\mathbb R}C_M(u)^*\mathcal N(u)C_M(u)\,du>0,
       &\mathscr R_M^N&=\operatorname{ran}\mathbf N_M.
 \end{aligned}
 \tag{FC.29}
\]
The coefficient adjoints in this subsection use the displayed literal
basis and its ordinary Euclidean coordinate inner product. The computed
metric pulled back from the normal image is \(H_M^N\). The original
source-amplitude metric and their exact comparison operator remain
\[
 \begin{gathered}
 H_M^0=\int_{\mathbb R}C_M(u)^*W(u)C_M(u)\,du>0,\qquad
 K_M^{N:0}=(H_M^0)^{-1}H_M^N,\\
 (K_M^{N:0})^*H_M^0=H_M^0K_M^{N:0}=H_M^N,\qquad
 p^*H_M^Nq=p^*H_M^0K_M^{N:0}q.
 \end{gathered}
 \tag{FC.29a}
\]
Indeed \(W>0\) and a nonzero vector polynomial \(C_Mp\)
give \(H_M^0>0\), with finite entries by the original moment
bounds. Direct multiplication proves the remaining identities, so
\(K_M^{N:0}\) is self-adjoint and strictly positive in the
\(H_M^0\) inner product. This is the exact operator between
the two retained forms; no identity of their norms is assumed.
All columns of \(\mathbf N_M\) and \(\mathbf T_M\)
belong to \(\mathscr K\). For the first map, projection contraction
gives \(\|N_uC_Mp\|\le\|j_u'C_Mp\|\), whose square is
integrable by the polynomially weighted Schwartz estimates.
For the second, it is the projection by \(\Pi_u\) of
\(\partial_u(j_uC_Mp)\), and the same estimates apply.
Equation (FC.26) proves strict positivity of \(H_M^N\) for every
nonzero coefficient vector. In particular the normal map is injective
on the whole original finite source, including nonzero polynomials
whose original arithmetic class is a supported zero.

\begin{theorem}[The actual normal image recovers the source and both graph parts]
The inverse on \(\mathscr R_M^N\), and its precise ambient formula,
are
\[
 \begin{gathered}
 L_M^N=(H_M^N)^{-1}\mathbf N_M^*:\mathscr K\to\mathbb C^{q_M},
 \qquad
 \mathbf N_M^*r=\int_{\mathbb R}C_M(u)^*N_u^*r(u)\,du,\\
 L_M^N\mathbf N_M=I_{q_M},\qquad
 \Pi_M^N=\mathbf N_ML_M^N
      \text{ is the orthogonal projection onto }\mathscr R_M^N.
 \end{gathered}
 \tag{FC.30}
\]
The full original derivative graph is
\(\mathscr G_M=\operatorname{ran}(\mathbf T_M,\mathbf N_M)\).
It is recovered from its normal image by the exact linear bijection
\[
 \begin{aligned}
 \mathsf{Rec}_M:\mathscr R_M^N&\longrightarrow\mathscr G_M,
       &r&\longmapsto(\mathbf T_ML_M^Nr,r),\\
 \mathsf{Rec}_M^{-1}(t,r)&=r,\qquad
 &p&=L_M^Nr,\quad P=P_{L_M^Nr}.
 \end{aligned}
 \tag{FC.31}
\]
Its complete metric identity retains both components:
\[
 \begin{gathered}
 H_M^{\partial}=\mathbf T_M^*\mathbf T_M+H_M^N,\qquad
 \|(\mathbf T_Mp,\mathbf N_Mp)\|^2=p^*H_M^{\partial}p,\\
 \|\mathsf{Rec}_Mr\|^2
      =\|r\|^2+\|\mathbf T_ML_M^Nr\|^2
      =(L_M^Nr)^*H_M^{\partial}(L_M^Nr).
 \end{gathered}
 \tag{FC.32}
\]
In particular reconstruction is not asserted to preserve the normal
norm alone; (FC.32) specifies its full metric.
\end{theorem}

\begin{proof}
Every entry of the integral in (FC.30) converges by
Cauchy--Schwarz, since it pairs \(r\) with one of the finitely
many \(L^2\) columns of \(\mathbf N_M\). Pairing against an
arbitrary coefficient vector proves that this integral is its
Hilbert adjoint. Therefore
\(L_M^N\mathbf N_M=(H_M^N)^{-1}H_M^N=I\).
Multiplication shows that \(\Pi_M^N\) is self-adjoint and
idempotent, that its image is \(\mathscr R_M^N\), and that it
fixes every vector of that image. This proves the projection and
inverse claims. The image is finite dimensional and hence closed;
all displayed maps are bounded.
If \(r=\mathbf N_Mp\), (FC.30) makes its reconstructed pair
exactly \((\mathbf T_Mp,\mathbf N_Mp)\). Conversely every pair
in \(\mathscr G_M\) has this form. Injectivity of
\(\mathbf N_M\) gives both uniqueness of \(p\) and the inverse
in (FC.31). The two graph components are pointwise orthogonal by
\(j_u^*N_u=0\); their squared norms give (FC.32).
Their addition consequently is the original derivative amplitude
and has the same full graph norm, as already proved in (FC.7).
\end{proof}

The full arithmetic derivative observation can now be specified on
this actual normal image. Put
\[
 \begin{gathered}
 I=(h(s_1),\ldots,h(s_k)),\qquad
 E=\mathbb C[\mathbf s]/I=E_h^{\otimes k},\qquad
 U(\mathbf s)=\prod_i v_h(s_i),\\
 \upsilon^{(k)}=j_IU\in E^\times,\qquad
 \partial_S^{\rm rel}=\frac1k\sum_i\partial_{s_i},\qquad
 \beta_h=(\upsilon^{(k)})^{-1}j_I(\partial_S^{\rm rel}U).
 \end{gathered}
 \tag{FC.33}
\]
Here \(j_I\) means the complete polynomial Hermite class in the
original product remainder coordinates. It retains every order in
each factor. The entire function \(U\) is a unit at the finite
packet; its inverse and every derivative used in \(\beta_h\)
are taken in those full local algebras. Thus this formula uses no
global division across a zero of \(v_h\). If \(h=1\), the
target \(E\) is the zero algebra and all arithmetic observations
below are its unique zero map; the preceding analytic reconstruction
continues unchanged.

For a nonempty packet define the finite matrix
\(\mathsf J_M^{\log}:\mathbb C^{q_M}\to E\) by its original
source columns. Mellin differentiation and (FC.9) give their exact
formulas
\[
 \begin{aligned}
 (\mathsf J_M^{\log})_{\alpha,n}
 &=J^{(k)}\mathscr L_k\mathcal T_h^{(k)}P_{\alpha,n}
   =j_I\bigl(\partial_S^{\rm rel}(UP_{\alpha,n})\bigr)\\
 &=n\upsilon^{(k)}[\theta_\alpha S^{n-1}]_I
       +\beta_h\upsilon^{(k)}[\theta_\alpha S^n]_I.
 \end{aligned}
 \tag{FC.34}
\]
For \(n=0\) the first term is defined to be zero, and no negative
power is used. To verify the last equality, the original coordinate
map gives \(\partial_S^{\rm rel}S=1\) and
\(\partial_S^{\rm rel}z_i=0\). Hence the relative polynomial
does not acquire a derivative term, while differentiating \(S^n\)
gives precisely its original factor \(n\). The other product-rule
term is the full \(\beta_h\) just defined. Equivalently
\(\beta_h=k^{-1}\sum_i j_I(v_h'/v_h)(s_i)\), with the
quotient taken only in its actual unit neighbourhoods. Both factors
\(g=2\xi\) and every Taylor coefficient remain in these columns.

The exact observation and the complete commuting source path are
\[
 \begin{gathered}
 \mathsf{Obs}_M^{\log}:\mathscr R_M^N\to E,\qquad
 \mathsf{Obs}_M^{\log}(r)
       =\mathsf J_M^{\log}(H_M^N)^{-1}\mathbf N_M^*r,\\
 \mathsf{Obs}_M^{\log}(\mathbf N_Mp)
       =J^{(k)}\mathscr L_k\mathcal T_h^{(k)}P_p,\\
 \mathscr R_M^N\xrightarrow{\mathsf{Rec}_M}\mathscr G_M
 \xrightarrow{\operatorname{Add}}
     i\mathscr U_k\mathscr L_k\mathcal T_h^{(k)}\mathcal P_M^\theta,\\
 i\mathscr U_k\mathscr L_k\mathcal T_h^{(k)}\mathcal P_M^\theta
 \xrightarrow{-i\mathscr U_k^{-1}}
     \mathscr L_k\mathcal T_h^{(k)}\mathcal P_M^\theta
 \xrightarrow{J^{(k)}}E.
 \end{gathered}
 \tag{FC.35}
\]
Substitute \(r=\mathbf N_Mp\) in (FC.30)--(FC.31).
Graph addition gives
\(\partial_u\mathscr U_k\mathcal T_h^{(k)}P_p
=i\mathscr U_k\mathscr L_k\mathcal T_h^{(k)}P_p\),
and the last two arrows give exactly the columns in (FC.34).
This proves equality of the two descriptions of the observation.
All functions in the penultimate space are actual finite tensor test
functions, since logarithmic multiplication preserves \(\mathscr B\).
Consequently this is an arithmetic observation on the specified
finite normal image with a proved source inverse. The ambient formula
for \(L_M^N\) also defines an orthogonal projection, but no
arithmetic-jet interpretation of an arbitrary unrelated measurable
normal section is required by (FC.35).

The reconstruction is unchanged under a declared smooth invertible
relative frame. In that frame its coefficient matrix is
\(C_M^C=C(u)^{-1}C_M\), and the exact identities in (FC.8) give
\[
 \begin{gathered}
 N^CC_M^C=NC_M,\qquad
 j^C\bigl((C_M^C)'+\Gamma^CC_M^C\bigr)
                  =j(C_M'+\Gamma C_M),\\
 (H_M^N)^C=H_M^N,\qquad
 (\mathbf N_M)^C=\mathbf N_M,\qquad
 (\mathbf T_M)^C=\mathbf T_M.
 \end{gathered}
 \tag{FC.36}
\]
The derivative of \(C^{-1}\) cancels exactly against the
\(C^{-1}C'\) connection term. Thus the original normal-image
inverse, the full reconstructed graph and the arithmetic observation
are the same maps on the same source after its explicit coefficient
transport. Their fixed-support linear lifts retain the original
label and all supported zeros. In particular a nonzero polynomial
relation can have a nonzero normal image whose reconstructed
arithmetic derivative is the full conormal response (FC.34);
the source relation is retained before either quotient observation.
